% =============================================================================
% Chapter 6 — Conclusion and Research Agenda
% =============================================================================

\chapter{Conclusion and Research Agenda}
\label{ch:conclusion}

\section{What the Dissertation Set Out to Do}
\label{sec:set-out}

This dissertation began with an observation and a wager. The observation: the same data infrastructures that promise to optimise the civilian city — its traffic, its hospitals, its policing — are structurally available to the apparatus that destroys cities. The wager: that this is not a metaphor but a mechanism, with identifiable economic, institutional, and epistemic components, and that a systematic review of three and a half decades of literature could test it. Chapter~2 assembled the mechanism from five theoretical elements: Arrow's information-good economics, Deleuze's control societies, De Landa's targetable city, Michael's epistemic authority, and the urbicide canon of King, Graham, and Weizman. Chapters~3 and~4 built the test: a PRISMA-guided review of 2,713 records screened across two identification streams, yielding a coded corpus of 270, mapped against the framework's four themes.

\section{What It Found}
\label{sec:what-it-found}

Three findings stand. First, the framework's elements are each confirmed, but in segregated literatures. Dual-use transfer is documented exhaustively in both directions — yet its economic grammar is untheorised, leaving Arrow's mechanism as the dissertation's first contribution. The programmable-city ontology is common ground for smart-city advocates, their critics, and military doctrine alike. Epistemic authority is shown to be institutionally manufactured and constitutively racialised, extending Michael beyond his original formulation. And the destruction canon is thriving, cyclical, and increasingly quantified.

Second, the intersection is empty. Only five records substantively address algorithmically mediated spatial destruction, all since 2021. The literature of how algorithms govern and the literature of how cities are destroyed run in parallel — a gap reproduced in the review's own architecture, where keyword search and citation chaining recovered the two sides independently. The dissertation's central empirical contribution is this map of a junction that has not yet been built.

Third, dual-use operates as a warrant, not merely a flow. The category through which export law restrains technology transfer is the same category through which targeting doctrine licenses strikes on civilian fabric — ``hospitals, schools, mosques\dots\ struck under the rationale of dual-use'' \parencite{benguita2025}. The collapse of the civilian/military distinction is the targeter's resource as much as the critic's concern. This finding, unanticipated in the framework, is the review's sharpest policy-relevant result.

\section{An Agenda for Research}
\label{sec:agenda}

The empty intersection defines the agenda. Five directions follow.

\textbf{1. Name the field.} Algorithmically mediated urbicide needs its own literature before it can have its own debates. The five-record beachhead \parencite{hourani2024,lotfijam2025,zavrsnik2021,graham2008,holail2024} should be consolidated — through a dedicated special issue or state-of-the-field review — into a citable research programme connecting critical geography, surveillance studies, and International Humanitarian Law (IHL) scholarship.

\textbf{2. Quantify the junction.} The corpus is 7\% quantitative. Remote-sensing damage assessment \parencite{holail2024} shows that the sensing stack's documentary use is feasible; systematic work pairing targeting disclosures with damage signatures could move the field from assertion to measurement.

\textbf{3. Extend the comparative base.} The destruction literature's Palestine concentration is both its strength and its constraint. Kurdish Turkey, Belfast, Rio, Kigali, and Basrah appear in the corpus as comparative footholds; systematic cross-case work would test whether the cyclical model — renewal, razing, reconstruction — generalises, and whether the algorithmic mediation of destruction differs across regime types.

\textbf{4. Institutionalise counter-forensics.} The corpus's clearest existing model of contestation is the counter-forensic use of the same stacks that target (Forensic Architecture; \cite{holail2024}). What mandates, funding lines, and evidentiary standards would make that practice a durable check rather than an activist art? This is the agenda's most direct policy question.

\textbf{5. Follow the ontologies into the era of autonomous weapons.} The review showed transfer moving through data models and interoperability standards rather than discrete artefacts. This finding connects directly to the urgent field of AI interpretability and alignment: the target ontologies, urban data models, and ``pattern of life'' schemas that travel between civilian platforms and military systems encode assumptions about what constitutes a legitimate target, a normal civilian, and an anomalous pattern. As lethal autonomous weapon systems (LAWS) transition from development to deployment, these ontologies become the de facto rules of engagement — machine-readable categories that determine who lives and who dies, with no human interpreter in the loop. Research is needed on the specific ontologies that mediate this transition, on the vendor contracts that carry them into operational contexts, and on the alignment failures they embed. This is empirical, document-based work that public administration is well placed to do, and it is the agenda's most consequential item.

\section{Closing}
\label{sec:closing}

The introduction opened with a loop: Palantir's software moving from insurgent-mapping to NHS bed management, and urban-analytic sensing moving from peacetime planning to target generation in Gaza. The systematic map shows that loop is not an anecdote but a structure — documented on both sides, theorised on neither. If the review has a single lesson for policy, it is that the relevant governance question is no longer whether a given urban technology can be turned to destruction; the corpus answers that in the affirmative, eighty-six times over. The question is whether democratic institutions will develop the categories — ontological, legal, evidentiary — to see the turn before it completes. At present, as the five records at the empty centre of this map attest, the literature has barely begun to look. That is the gap this dissertation has mapped, and the work it invites.